\input{preamble.tex}
\title{Отчет по лабораторной работе № 5. \\ Простые сети в GNS3. Анализ трафика}
\author{Данила Стариков \\ НПИбд-02-22}
\date{2024}

\begin{document}

\maketitle
\newpage

\tableofcontents

\newpage
\section{Цель работы}
Построение простейших моделей сети на базе коммутатора и маршрутизаторов FRR и VyOS в GNS3, анализ трафика посредством Wireshark.

\newpage
\section{Выполнение работы}
\subsection{Моделирование простейшей сети на базе коммутатора в GNS3}
\begin{enumerate}
\item Запустили GNS3 VM и GNS3. Создали новый проект.
\item В рабочей области GNS3 разместили коммутатор Ethernet и два VPCS (рис. \ref{01}). Изменили название устройства, включив в имя устройства имя учётной записи. Коммутатору присвоили название \texttt{msk-dastarikov-sw-01}. Соединили VPCS с коммутатором. Отобразили обозначение интерфейсов соединения.

\begin{center}
    \centering
    \includegraphics[width=\textwidth]{../images/image01.png}
    \captionof{figure}{Топология сети.}
    \label{01}
\end{center}


\item Задали IP-адреса VPCS. Для этого с помощью меню, вызываемого правой кнопкой мыши, запустили Start, например, PC-1, затем вызвали его терминал Console. Для просмотра синтаксиса возможных для ввода команд набрали {\tt /?}.
Для задания IP-адреса 192.168.1.11 в сети 192.168.1.0/24 ввели (рис. \ref{02}):
\begin{minted}{bash}
ip 192.168.1.11/24 192.168.1.1
\end{minted}

\item  Аналогичным образом задали IP-адрес 192.168.1.12 для PC-2 (рис. \ref{03}):

\begin{minted}{bash}
ip 192.168.1.12/24 192.168.1.1
\end{minted}

  
\begin{center}
    \centering
    \includegraphics[width=\textwidth]{../images/image02.png}
    \captionof{figure}{Задание IP-адреса для PC-1.}
    \label{02}
\end{center}

\begin{center}
    \centering
    \includegraphics[width=\textwidth]{../images/image03.png}
    \captionof{figure}{Задание IP-адреса для PC-2.}
    \label{03}
\end{center}


\item Проверили работоспособность соединения между PC-1 и PC-2 с помощью
команды \texttt{ping} (рис. \ref{04} и \ref{05}).

  \begin{center}
    \centering
    \includegraphics[width=\textwidth]{../images/image04.png}
    \captionof{figure}{Отправка эхо-запроса от PC-1 к PC-2.}
    \label{04}
\end{center}

\begin{center}
    \centering
    \includegraphics[width=\textwidth]{../images/image05.png}
    \captionof{figure}{Отправка эхо-запроса от PC-2 к PC-1.}
    \label{05}
\end{center}


\item Остановили в проекте все узлы (меню GNS3 \texttt{Control} \(\rightarrow\) \texttt{Stop all nodes}).
\end{enumerate}

\subsection{Анализ трафика в GNS3 посредством Wireshark}
\begin{enumerate}
\item Запустили на соединении между PC-1 и коммутатором анализатор трафика. Для этого щёлкнули правой кнопкой мыши на соединении, выбрали в меню \texttt{Start capture}. % при необходимости можете скорректировать название DUMP-файла
Запустился Wireshark, а в проекте GNS3 на соединении появится значок лупы.
\item В проекте GNS3 стартовали все узлы (меню GNS3 \texttt{Control} \(\rightarrow\) \texttt{xStart/Resume all nodes}). В окне Wireshark отобразилась информация по протоколу ARP (рис. \ref{06}).

\begin{center}
    \centering
    \includegraphics[width=\textwidth]{../images/image06.png}
    \captionof{figure}{Просмотр перехваченных пакетов по протоколу ARP.}
    \label{06}
\end{center}

Оба устройства отправляют ARP-запросы для определения MAC-адресов устройств, чей IP-адрес им уже известен.

\item В терминале PC-2 посмотрели информацию по опциям команды \texttt{ping}, введя \texttt{ping /?}. Затем сделали один эхо-запрос в ICMP-моде к узлу PC-1 (рис. \ref{07}).
  
  
\begin{center}
    \centering
    \includegraphics[width=\textwidth]{../images/image07.png}
    \captionof{figure}{Перехваченный пакет эхо-запроса в ICMP-моде.}
    \label{07}
\end{center}



  % В окне Wireshark проанализируйте полученную информацию, дайте пояснения в отчёте.
\item Сделали один эхо-запрос в UDP-моде к узлу PC-1 (рис. \ref{08} и \ref{09}).
    
\begin{center}
    \centering
    \includegraphics[width=\textwidth]{../images/image08.png}
    \captionof{figure}{Отправка эхо-запроса в UDP-моде.}
    \label{08}
\end{center}

\begin{center}
    \centering
    \includegraphics[width=\textwidth]{../images/image09.png}
    \captionof{figure}{Перехваченный пакет эхо-запроса в UDP-моде.}
    \label{09}
\end{center}

\item Сделали один эхо-запрос в TCP-моде к узлу PC-1 (рис. \ref{10} и \ref{11}).
  
\begin{center}
    \centering
    \includegraphics[width=\textwidth]{../images/image10.png}
    \captionof{figure}{Отправка эхо-запроса в TCP-моде.}
    \label{10}
\end{center}

\begin{center}
    \centering
    \includegraphics[width=\textwidth]{../images/image11.png}
    \captionof{figure}{Перехваченный пакет эхо-запроса в TCP-моде.}
    \label{11}
\end{center}

Для отправки TCP-пакета сначала выполняется процедура трехстороннего рукопожатия (\texttt{TCP handshake}), которое подробнее разбиралось в лабораторной работе №3. После установления соединения отправляется сам эхо-запрос.

\item Остановили захват пакетов в Wireshark.
\end{enumerate}

\subsection{Моделирование простейшей сети на базе маршрутизатора VyOS в GNS3}
\begin{enumerate}
\item Запустили GNS3 VM и GNS3. Создали новый проект.
\item В рабочей области GNS3 разместили VPCS, коммутатор Ethernet и маршрутизатор VyOS (рис. \ref{21}).
  
\begin{center}
    \centering
    \includegraphics[width=\textwidth]{../images/image21.png}
    \captionof{figure}{Топология сети.}
    \label{21}
\end{center}

\item Изменили отображаемые названия устройств. Коммутатору присвоили название по принципу \texttt{msk-dastarikov-sw-0x}, маршрутизатору — по принципу \texttt{msk-dastarikov-gw-0x}, VPCS — по принципу \texttt{PCx-dastarikov}, вместо x — порядковый номер устройства.
\item Включили захват трафика на соединении между коммутатором и маршрутизатором.
\item Запустили все устройства проекта. Открыли консоль всех устройств проекта.
\item Настроили IP-адресацию для интерфейса узла PC1 (рис. \ref{12}):
  \begin{minted}{bash}
    ip 192.168.1.10/24 192.168.1.1
    save
    show ip
  \end{minted}
  
\begin{center}
    \centering
    \includegraphics[width=\textwidth]{../images/image12.png}
    \captionof{figure}{Задание IP-адреса для PC-1.}
    \label{12}
\end{center}


\item Настроили маршрутизатор VyOS (рис. \ref{14}-\ref{16}):
    \begin{enumerate}
\item После загрузки ввели логин vyos и пароль vyos:
  \begin{minted}{bash}
    vyos login: vyos
    Password:
  \end{minted}

\item Перешли в режим конфигурирования:
  \begin{minted}{bash}
    vyos@vyos$ configure
    vyos@vyos#
  \end{minted}

\item Изменили имя устройства (вместо dastarikov укажите свою учётную запись):
  \begin{minted}{bash}
    vyos@vyos#set system host-name msk-dastarikov-gw-01
  \end{minted}
  
\begin{center}
    \centering
    \includegraphics[width=\textwidth]{../images/image14.png}
    \captionof{figure}{Изменение имени устройства.}
    \label{14}
\end{center}

\item Задали IP-адрес на интерфейсе eth0:
  \begin{minted}{bash}
    vyos@vyos# set interfaces ethernet eth0 address 192.168.1.1/24
  \end{minted}
  
\begin{center}
    \centering
    \includegraphics[width=\textwidth]{../images/image15.png}
    \captionof{figure}{Задание IP-адреса на интерфейсе eth0.}
    \label{15}
\end{center}


\item Посмотрели внесённые в конфигурацию изменения:
  \begin{minted}{bash}
    vyos@vyos# compare
  \end{minted}
\item Применили изменения в конфигурации и сохранили саму конфигурацию:
  \begin{minted}{bash}
    vyos@vyos# commit
    vyos@vyos# save
  \end{minted}
  
\begin{center}
    \centering
    \includegraphics[width=\textwidth]{../images/image17.png}
    \captionof{figure}{Сохранение изменений.}
    \label{17}
\end{center}


\item Посмотрили информацию об интерфейсах маршрутизатора:
  \begin{minted}{bash}
    vyos@vyos# show interfaces
  \end{minted}
  
  \begin{center}
    \centering
    \includegraphics[width=\textwidth]{../images/image16.png}
    \captionof{figure}{Просмотр информации об интерфейсах маршрутизатора.}
    \label{16}
  \end{center}

\item Вышли из режима конфигурирования:
  \begin{minted}{bash}
    vyos@vyos# exit
    vyos@vyos$
  \end{minted}
    \end{enumerate}
\item Проверили подключение. Узел PC1 должен успешно отправлять эхо-запросы на адрес маршрутизатора 192.168.1.1  (рис. \ref{18} и \ref{19}).

\begin{center}
    \centering
    \includegraphics[width=\textwidth]{../images/image18.png}
    \captionof{figure}{Отправка эхо-запроса от PC1 к маршрутизатору.}
    \label{18}
\end{center}
  
\begin{center}
    \centering
    \includegraphics[width=\textwidth]{../images/image19.png}
    \captionof{figure}{Просмотр перехваченных пакетов.}
    \label{19}
\end{center}

\item Остановили захват пакетов в Wireshark. Остановили все устройства в проекте. Завершили работу с GNS3.
\end{enumerate}

\section{Выводы}
В результате выполнения лабораторной работы построили простейшие модели сети на базе коммутатора и маршрутизатора VyOS в GNS3 и провели анализ трафика посредством Wireshark.
\end{document}
